\section{Resultados}
~\label{sec:resultados}

La validación del sistema se realizó en un entorno controlado 
montado a partir de una red local de catorce nodos: doce endpoints 
con el sensor de captura instalado, un servidor central que aloja 
el pipeline de análisis, la API y la base de datos 
contenedorizada, y un dispositivo móvil con la aplicación 
conectado a la red privada Tailscale. La evaluación se organizó en 
torno a tres ejes que responden de forma directa a los objetivos 
específicos del proyecto: el rendimiento del sistema bajo carga 
operativa (Sección~\ref{subsec:res-rendimiento}), el comportamiento 
del motor de detección de anomalías frente a eventos inducidos 
(Sección~\ref{subsec:res-motor}), y la verificación funcional 
extremo a extremo desde los sensores hasta la aplicación móvil 
(Sección~\ref{subsec:res-e2e}).

\subsection{Rendimiento del sistema bajo carga operativa}
~\label{subsec:res-rendimiento}

La primera serie de mediciones buscó caracterizar la huella del 
sistema en régimen de operación normal. Se instrumentaron los tres 
componentes principales durante diez minutos de operación 
continuada con los doce sensores conectados simultáneamente y con 
ciclos de análisis activos, registrando muestras cada cinco 
segundos sobre CPU, memoria residente, hilos activos y tráfico de 
red. La Tabla~\ref{tab:rendimiento} consolida los resultados 
principales.

\begin{table}[ht]
\centering
\caption{Consumo de recursos durante la validación del sistema. Se reportan valores promedio y máximo observados sobre 120 muestras por componente.}
~\label{tab:rendimiento}
\small
\begin{tabular}{@{}lrr@{}}
\toprule
\textbf{Componente} & \textbf{Promedio} & \textbf{Máximo} \\
\midrule
\multicolumn{3}{l}{\textit{Sensor de captura (por endpoint)}} \\
\quad CPU            & 2{,}43 \% & 2{,}70 \% \\
\quad RAM residente  & 94{,}54 MB & 98{,}68 MB \\
\quad Hilos activos  & 4 & 4 \\
\quad Tráfico saliente & 4{,}62 KB/s & --- \\
\midrule
\multicolumn{3}{l}{\textit{Motor de análisis (servidor central)}} \\
\quad CPU            & 7{,}9 \% & 9{,}1 \% \\
\quad RAM residente  & 156{,}3 MB & 184{,}2 MB \\
\quad Hilos activos  & 22 & 28 \\
\midrule
\multicolumn{3}{l}{\textit{API REST y WebSocket (servidor central)}} \\
\quad CPU            & 0{,}20 \% & 0{,}21 \% \\
\quad RAM residente  & 70{,}1 MB & 88{,}2 MB \\
\quad Hilos activos  & 1 & 1 \\
\midrule
\multicolumn{3}{l}{\textit{Base de datos PostgreSQL (contenedor)}} \\
\quad CPU            & 1{,}4 \% & 3{,}1 \% \\
\quad RAM residente  & 58{,}7 MB & 64{,}1 MB \\
\quad Escritura disco (10 min) & --- & 95 MB \\
\bottomrule
\end{tabular}
\end{table}

Tres observaciones emergen de estos valores. La primera concierne 
al sensor de captura: su huella permanece por debajo de 100~MB de 
memoria residente y por debajo del 3~\% de CPU en el endpoint, lo 
que satisface el requerimiento no funcional de bajo impacto que 
condicionó el diseño del módulo de captura y la decisión de 
sanitizar el tráfico en origen. La constancia del número de hilos 
a lo largo de las 120~muestras confirma que el modelo 
productor-consumidor descrito en la Sección~\ref{subsec:sensor} 
opera sin fugas ni hilos huérfanos.

La segunda observación refiere a la eficiencia de la sanitización 
dinámica. Durante los diez minutos de medición sobre un endpoint 
representativo, el sensor observó 7{,}60~MB de tráfico entrante a 
la interfaz de red y transmitió 2{,}71~MB al servidor central, lo 
que corresponde a una reducción del volumen transmitido del 
64~\% respecto al observado. Esta cifra confirma 
cuantitativamente que el recorte a 96~bytes para tráfico genérico 
y 300~bytes para DNS y TLS preserva la información diagnóstica 
necesaria al tiempo que descarta los payloads sensibles, tal como 
se justificó en el marco legal de la Sección~\ref{sec:marco-legal}.

La tercera observación concierne al margen de recursos del 
servidor central. La carga combinada del motor, la API y la base 
de datos se mantuvo por debajo del 10~\% de CPU y no superó 
350~MB de memoria combinada en ningún momento de la medición. 
Este margen sustancial respalda la decisión arquitectónica de 
separar el pipeline de procesamiento de la capa de exposición como 
procesos independientes comunicados mediante el patrón Shared 
Database: ambos procesos disponen de recursos holgados para 
absorber picos de carga sin interferirse mutuamente.

\subsection{Comportamiento del motor de detección de anomalías}
~\label{subsec:res-motor}

[TODO — esta subsección requiere:
  (a) inducir una anomalía real en el sistema durante su operación,
  (b) registrar el timestamp de inducción y el de detección,
  (c) capturar el valor de |Z| reportado,
  (d) tomar screenshot de la alerta recibida en la app.
Redactar esta subsección una vez se ejecute el experimento.]

\subsection{Verificación funcional extremo a extremo}
~\label{subsec:res-e2e}

[TODO — esta subsección requiere:
  (a) screenshots de la app móvil: login, dashboard, inventario,
      detalle de dispositivo, lista de alertas, notificación push,
  (b) captura del sistema operando desde una red distinta a la LAN
      monitoreada (demostrando acceso remoto vía Tailscale),
  (c) opcionalmente, una secuencia de capturas mostrando el flujo
      completo desde el evento hasta la notificación.
Redactar esta subsección una vez se tomen las capturas.]